\begin{Theorem}
\label{thm:trace-bounded-right-inverse}
设 $s>1/2$. 那么存在有界线性映射 $Z:H^{s-1/2}(\mathbb{R}^{m-1})\to H^s(\mathbb{R}^m)$, 使 $TZ$ 为恒等映射.
\end{Theorem}

\begingroup
\postdisplaypenalty=10000
\renewcommand{\endProof}{\par\nobreak\hfill$\blacksquare$\par\vspace{0.1cm}}
\begin{Proof}
将用 Fourier 变换显式构造 $Z$. 由稠密性, 只需对 $\phi\in\mathcal{D}(\mathbb{R}^{m-1})$ 定义 $Z\phi$; 随后可由连续性延拓. 令
\begin{equation}\label{eq:trace-right-inverse-explicit-fourier-extension}
Z\phi:=u(\bm{x})=\dfrac{1}{(2\pi)^{(m-1)/2}K_s}\int_{\mathbb{R}^m}e^{i\bm{\xi}\cdot\bm{x}}\dfrac{(1+|\bm{\xi}^{\prime}|^2)^{s-1/2}}{(1+|\bm{\xi}|^2)^s}\hat\phi(\bm{\xi}^{\prime})\,\mathrm{d}\bm{\xi},
\end{equation}
其中令
\begin{equation}\label{eq:trace-right-inverse-normalizing-integral}
K_s=\int_{-\infty}^{\infty}(1+y^2)^{-s}\,\mathrm{d}y.
\end{equation}
若 $x_m=0$, 可完成关于 $\xi_m$ 的积分, 得到
\begin{equation}\label{eq:trace-right-inverse-hyperplane-restriction}
u(\bm{x}^{\prime},0)=(2\pi)^{-(m-1)/2}\int_{\mathbb{R}^{m-1}}e^{i\bm{\xi}^{\prime}\cdot\bm{x}^{\prime}}\hat\phi(\bm{\xi}^{\prime})\,\mathrm{d}\bm{\xi}^{\prime}=\phi(\bm{x}^{\prime}).
\end{equation}
这表明 $TZ\phi=\phi$. 还需证明 $Z$ 的连续性. 有
\begin{equation}\label{eq:trace-right-inverse-weighted-norm-bound}
\begin{aligned}
\|u\|_{s,2}^2&\leq C_1\int_{\mathbb{R}^m}|\hat u(\bm{\xi})|^2(1+|\bm{\xi}|^2)^s\,\mathrm{d}\bm{\xi}\\
&=C_2\int_{\mathbb{R}^{m-1}}(1+|\bm{\xi}^{\prime}|^2)^{2s-1}|\hat\phi(\bm{\xi}^{\prime})|^2\int_{-\infty}^{\infty}(1+|\bm{\xi}|^2)^{-s}\,\mathrm{d}\xi_m\,\mathrm{d}\bm{\xi}^{\prime}\\
&\leq C_3\int_{\mathbb{R}^{m-1}}(1+|\bm{\xi}^{\prime}|^2)^{s-1/2}|\hat\phi(\bm{\xi}^{\prime})|^2\,\mathrm{d}\bm{\xi}^{\prime}.
\end{aligned}
\end{equation}
证明完毕.
\end{Proof}
\endgroup
